\documentclass[11pt]{article}
\usepackage[margin=1in]{geometry}
\usepackage{amsmath,amssymb}
\usepackage[T1]{fontenc}
\usepackage[utf8]{inputenc}
\usepackage{lmodern}
\usepackage{microtype}
\usepackage{booktabs}
\usepackage{xurl}
\usepackage[hidelinks]{hyperref}
\DeclareUnicodeCharacter{220E}{\ensuremath{\square}}
\setlength{\emergencystretch}{3em}
\title{Restore-Sensitive Support Normal Forms for Multi-Tag Quantum Compilation}
\author{Sze Chun Yiu}
\date{}
\begin{document}
\maketitle

\begin{abstract}
Support bounds in exact quantum compilation often arise from a
restricted deletion proof rather than from an intrinsic lower bound of
the compiler. We isolate that distinction for a MultiTag-TARE grammar.
The combinatorial zero-sum invariant used in the proof is donor
mathematics: for a finite abelian signature group \texttt{H} and a fixed
allowed alphabet \texttt{A}, let \(\operatorname{zsf}(H;A)\) be the
longest zero-sum-free sequence over \texttt{A}. The new content is the
compiler-side contract that makes this donor invariant usable.

We state the deletion theorem with the quantifiers needed by a real
compiler: each generator has an instance-level alphabet fixed before
optimization; a legal deletion preserves feasibility of the whole
instance; and repeated deletions cannot increase any generator's
support. Starting from an optimum and applying admissible zero-sum
deletions until a global fixed point is reached gives one optimum
satisfying the \(\operatorname{zsf}\) ceiling simultaneously for every
constrained generator. In an elementary binary quotient this implies the
familiar rank ceiling, but that rank statement is explicitly
binary-specific and is not itself novel.

For the registered arbitrary-block MultiTag grammar, one frame-support
deletion changes one argument of one charged local \texttt{b}-way
Restore term. The exact one-argument sensitivity of that term is
\texttt{b-1}; therefore the objective-deletion cone is

\(mu \ge (b-1)t_R\).

Inside this cone, every admitted instance has an optimum with

\[\operatorname{support}(R) \le  \operatorname{zsf}(H_R;A_R) \le  \operatorname{rank}(H_R) \le  s+1\]

for every frame \texttt{R}, where \texttt{A\_R} is the fixed
instance-level signature alphabet and \(H_R=<A_R>\). The frozen one-Tag,
three-block R6M family is the sharp control: the certificate ceiling is
two and independent compiler evidence establishes intrinsic support
\(kappa_R6M=2\). We do not claim general MultiTag sharpness, necessity
outside the cone, or physical quantum advantage.
\end{abstract}

\noindent\textbf{Keywords:} quantum compilation; Pauli strings; support normal forms; zero-sum deletion; sparse optimization

\hypertarget{scientific-question-and-novelty-boundary}{%
\section{1. Scientific question and novelty
boundary}\label{scientific-question-and-novelty-boundary}}

The question is not whether Davenport-type zero-sum thresholds exist;
they do, and that mathematics is prior work. The question is what a
compiler must establish before such a threshold becomes a valid
simultaneous normal form under its semantics and objective.

The surviving paper-level contributions are therefore:

\begin{enumerate}
\def\labelenumi{\arabic{enumi}.}
\item
  an explicit whole-instance deletion contract and global fixed-point
  proof that yields simultaneous support ceilings;
\item
  the exact Restore sensitivity \texttt{b-1} and resulting objective
  cone for the stated MultiTag grammar;
\item
  the compiler-specific separation between certificate ceiling and
  intrinsic support, with R6M as a sharp control.
\end{enumerate}

The following are \textbf{not} claimed as novel: Davenport constants,
zero-sum-free sequence theory, the binary rank corollary, generic
proof-system-relative complexity, or sparse-support theory.

\hypertarget{fixed-alphabets-and-zero-sum-free-sequences}{%
\section{2. Fixed alphabets and zero-sum-free
sequences}\label{fixed-alphabets-and-zero-sum-free-sequences}}

Let \texttt{H} be a finite abelian group written additively and let
\texttt{A\ subseteq\ H} be a finite alphabet fixed independently of the
optimum being bounded. A \textbf{subsequence} means an arbitrary
sub-multiset of positions; it need not be contiguous.

A sequence over \texttt{A} is zero-sum-free if no nonempty subsequence
has sum zero. Define

\(\operatorname{zsf}(H;A) = max({0} \cup {|W| : W is zero-sum-free over A})\).

The explicit zero in the maximum handles degenerate alphabets. For the
full nonzero alphabet the object reduces to the standard small-Davenport
threshold; for restricted alphabets it is an alphabet-specific variant.
We use \(\operatorname{zsf}\) only as notation for this donor-owned
combinatorial object.

\hypertarget{simultaneous-deletion-theorem}{%
\section{3. Simultaneous deletion
theorem}\label{simultaneous-deletion-theorem}}

Fix one optimization instance. For each constrained generator
\texttt{R}, let \texttt{A\_R} be the set of signatures realizable by any
admissible local state of that instance, fixed before optimization, and
let every active coordinate of \texttt{R} carry
\texttt{v\_q\ in\ A\_R\ subseteq\ H\_R}.

Assume:

\begin{enumerate}
\def\labelenumi{\arabic{enumi}.}
\item
  \textbf{nonzero total:} every feasible constrained generator satisfies
  \(sum_q v_q \ne 0\);
\item
  \textbf{whole-instance deletion soundness:} deleting any nonempty
  zero-sum subsequence of coordinates of \texttt{R} preserves all
  constraints of the full instance, not only constraints local to
  \texttt{R};
\item
  \textbf{objective dominance:} the same deletion does not increase the
  objective;
\item
  \textbf{support monotonicity:} deleting coordinates of one generator
  does not activate coordinates of another generator or otherwise
  increase any generator's support.
\end{enumerate}

\textbf{Theorem 1 (simultaneous alphabet ceiling).} Every admitted
instance has an exact optimum such that, simultaneously for every
constrained generator \texttt{R},

\(\operatorname{support}(R) \le \operatorname{zsf}(H_R;A_R)\).

\textbf{Proof.} Start from any optimum. If some generator \texttt{R} has
a signature sequence longer than \(\operatorname{zsf}(H_R;A_R)\), it
contains a nonempty zero-sum subsequence. Since the total signature of
\texttt{R} is nonzero, that subsequence cannot be the whole active
sequence. Delete it. Assumptions 2 and 3 preserve feasibility and
optimality, while total support strictly decreases and no other support
increases by Assumption 4. Repeat whenever any generator remains
reducible. Total support is a nonnegative integer, so the process
terminates at a global fixed point. At that fixed point no generator
contains a nonempty zero-sum subsequence; hence every generator's active
signature sequence is zero-sum-free over its fixed alphabet and has
length at most \(\operatorname{zsf}(H_R;A_R)\). This one terminal
optimum satisfies all generator bounds simultaneously. ∎

This global fixed-point argument is the required quantifier step; a
one-pass instruction to ``repeat for every generator'' is not sufficient
when generators may share constraints.

\hypertarget{binary-rank-is-only-a-corollary}{%
\section{4. Binary rank is only a
corollary}\label{binary-rank-is-only-a-corollary}}

If \(H=F_2^d\), any sequence of more than \texttt{d} vectors is linearly
dependent. The dependent positions form a nonempty zero-XOR subsequence,
so

\(\operatorname{zsf}(F_2^d;A) \le d\).

If \texttt{A} contains a basis, equality holds for the abstract deletion
language.

This statement is specific to elementary binary groups. It must not be
generalized to arbitrary finite abelian groups: for example, in
\texttt{Z\_n} with alphabet \texttt{\{1\}}, the sequence of \texttt{n-1}
ones is zero-sum-free even though the group has rank one.

A rank ceiling becomes an intrinsic compiler statement only after a
separate compiler lower witness.

\hypertarget{multitag-tare-grammar-and-the-nonzero-total-invariant}{%
\section{5. MultiTag-TARE grammar and the nonzero-total
invariant}\label{multitag-tare-grammar-and-the-nonzero-total-invariant}}

Fix \(b\ge 2\) ordered blocks and \(s\ge 0\) shared Tag Paulis
\texttt{S\_1,...,S\_s}. Each constrained frame Pauli \texttt{R} has an
anticommuting partner \texttt{R\textquotesingle{}}. At an active frame
coordinate \texttt{q}, define

\(v_q=(<R_q,R'_q>,<S_1q,R_q>,...,<S_sq,R_q>) in F_2^(s+1)\).

The XOR of the first components equals the global binary symplectic
product \texttt{\textless{}R,R\textquotesingle{}\textgreater{}}. Since
\texttt{R} and \texttt{R\textquotesingle{}} anticommute, \(<R,R'\ge 1\),
and therefore the total signature is nonzero. This discharges Theorem
1's first assumption for the registered grammar.

For each frame \texttt{R}, define \texttt{A\_R} at the instance level as
the set of signatures realizable by any admissible local frame state,
not the alphabet observed in a particular optimum. Set \(H_R=<A_R>\).

\hypertarget{exact-restore-sensitivity}{%
\section{6. Exact Restore sensitivity}\label{exact-restore-sensitivity}}

The local Restore term is

\[F_b(a_1,...,a_b)=1\]

when all \texttt{b} letters are the same nonidentity Pauli, and
otherwise equals the number of nonidentity letters.

\textbf{Lemma 2 (one-argument sensitivity).} Replacing one argument of
\texttt{F\_b} can increase its value by at most \texttt{b-1}, and the
bound is attained whenever the local Pauli alphabet contains at least
two distinct nonidentity letters.

\textbf{Proof.} Away from the all-equal nonidentity state, changing one
argument changes ordinary nonidentity Hamming weight by at most one.
Entering the special state lowers the value. Leaving the all-equal
nonidentity state by replacing one letter with a different nonidentity
letter changes the value from \texttt{1} to \texttt{b}, giving the exact
increase \texttt{b-1}. ∎

The registered MultiTag incidence contract assigns a frame-support
deletion to exactly one argument replacement in exactly one charged
local \texttt{F\_b} term. This incidence is part of the grammar; without
it the following objective cone would have to be multiplied by the
number of affected Restore terms.

\hypertarget{restore-sensitive-multitag-normal-form}{%
\section{7. Restore-sensitive MultiTag normal
form}\label{restore-sensitive-multitag-normal-form}}

The objective charges each active frame coordinate by at least
\texttt{mu}, includes arbitrary nonnegative Tag-support terms, and
charges the local Restore term by \(t_R\ge 0\).

Deleting \texttt{k} frame coordinates in a sound zero-sum subsequence
refunds at least \texttt{k\ mu}. By the incidence contract and Lemma 2
it adds at most \texttt{k(b-1)t\_R} in Restore cost. Hence deletion is
non-increasing whenever

\(mu \ge (b-1)t_R\).

\textbf{Theorem 3 (MultiTag support normal form).} Under the registered
grammar and incidence contract, if \(mu \ge (b-1)t_R\), every admitted
instance has an exact optimum satisfying

\[\operatorname{support}(R) \le  \operatorname{zsf}(H_R;A_R) \le  \operatorname{rank}(H_R) \le  s+1\]

simultaneously for every frame \texttt{R}.

\textbf{Proof.} Zero signature preserves the partner and Tag constraints
by the registered semantic map. The objective calculation above
establishes deletion dominance. Theorem 1 then gives the simultaneous
\(\operatorname{zsf}\) ceiling. Because \texttt{H\_R} is an elementary
binary subgroup of \texttt{F\_2\^{}(s+1)}, Section 4 gives the rank
bounds. ∎

The first inequality may be strictly stronger than rank when the fixed
admissible alphabet does not realize a basis of its generated subgroup.

\hypertarget{sharp-r6m-control}{%
\section{8. Sharp R6M control}\label{sharp-r6m-control}}

For the frozen one-Tag, three-block R6M family,

\(b=3\), \(s=1\), \(mu=2\), \(t_R=1\).

The objective lies on the cone boundary. The registered signature
alphabet realizes a basis of \texttt{F\_2\^{}2}, so the rank-only
deletion ceiling is two. The commit-bound parent evidence supplies the
all-size upper result and the independent support-one obstruction needed
to conclude

\(kappa_R6M=2\),

where \(kappa\) denotes the mathematical minimum uniform support bound
of the compiler family under the frozen objective. This exact value is a
compiler result; the zero-sum/rank argument alone would not establish
it.

No corresponding lower theorem is claimed for general MultiTag-TARE.

\hypertarget{relation-to-prior-work}{%
\section{9. Relation to prior work}\label{relation-to-prior-work}}

Zero-sum sequence theory owns Davenport-type thresholds and
restricted-alphabet variants. Sparse integer optimization owns general
support bounds for optimal solutions. Standard stabilizer formalism owns
the binary symplectic signature machinery. Proof-complexity and
formal-methods literatures own the general distinction between a proof
system and the object it certifies.

The residual contribution is narrower: the registered TARE/MultiTag
semantic signature, the exact Restore sensitivity and objective cone,
the whole-instance simultaneous deletion contract, and the sharp R6M
control.

\hypertarget{reproducibility-and-authority}{%
\section{10. Reproducibility and
authority}\label{reproducibility-and-authority}}

The R6M parent theorem, objective ledger and necessity witness are
commit-bound in the existing A1 evidence package. The finite verifiers
test the local Restore function and small-group signature mechanisms,
but all-size authority comes from the proofs above and from the cited
R6M parent theorem. Internal independent implementations are not
external replication.

\hypertarget{limitations}{%
\section{11. Limitations}\label{limitations}}

\begin{enumerate}
\def\labelenumi{\arabic{enumi}.}
\item
  The zero-sum object is donor mathematics; novelty is not claimed for
  \(\operatorname{zsf}\) itself.
\item
  The deletion theorem requires whole-instance soundness, objective
  dominance and support monotonicity.
\item
  The MultiTag cone uses the explicitly stated
  one-deletion/one-Restore-argument incidence contract.
\item
  Outside \(mu \ge (b-1)t_R\), the proof is silent; no larger-support
  necessity follows.
\item
  General MultiTag sharpness is open.
\item
  Structural support is not T count, depth, runtime, qubits, or quantum
  advantage.
\item
  Submission-date donor overlap and independent specialist proof review
  remain external scientific checks.
\end{enumerate}

\hypertarget{conclusion}{%
\section{12. Conclusion}\label{conclusion}}

The scientifically defensible result is not a new zero-sum invariant. It
is a compiler theorem specifying exactly when donor zero-sum mathematics
can be converted into a simultaneous support normal form. The global
fixed-point proof closes the multi-generator quantifier; the
instance-level alphabet removes solution-relative circularity; the
Restore incidence and exact \texttt{b-1} sensitivity expose the
objective condition; and R6M shows one case where the certificate
ceiling is genuinely intrinsic. That bounded result is the journal
object.

\hypertarget{tool-use-disclosure}{%
\section{Tool-use disclosure}\label{tool-use-disclosure}}

A generative language model assisted manuscript organization, language
revision, adversarial review, and submission-package preparation. The
listed author remains responsible for the mathematical statements,
proofs, references, executable claims, and final text.

\hypertarget{data-and-code-availability}{%
\section{Data and code availability}\label{data-and-code-availability}}

The source package accompanying this manuscript contains finite-group
and Restore-sensitivity control records used for implementation checks.
These files are provenance aids rather than theorem authority; the
displayed proofs and the cited parent theorem/witness carry the all-size
claims.

\hypertarget{author-contributions}{%
\section{Author contributions}\label{author-contributions}}

Sze Chun Yiu is the sole listed author and performed the authorship
contributions represented in this manuscript, including the formal
analysis, computational verification, and manuscript preparation.

\hypertarget{references}{%
\section{References}\label{references}}

\begin{enumerate}
\def\labelenumi{\arabic{enumi}.}
\item
  N. Schillo, A. Sturm, and R. Quay, ``TARE: Block Encoding Linear
  Combinations of Pauli Strings Without Ancilla State Preparation,''
  arXiv:2601.05740v4 {[}quant-ph{]} (2026).
\item
  I. Aliev, J. A. De Loera, F. Eisenbrand, T. Oertel, and R. Weismantel,
  ``The Support of Integer Optimal Solutions,'' \emph{SIAM Journal on
  Optimization} \textbf{28}, 2152--2157 (2018). DOI: 10.1137/17M1162792.
\item
  M. Freeze and W. A. Schmid, ``Remarks on a generalization of the
  Davenport constant,'' \emph{Discrete Mathematics} \textbf{310},
  3373--3389 (2010). DOI: 10.1016/j.disc.2010.07.028.
\item
  G. Wang, ``The universal zero-sum invariant and weighted zero-sum for
  infinite abelian groups,'' \emph{Communications in Algebra}
  \textbf{53}(4), 1581--1599 (2025). DOI: 10.1080/00927872.2024.2418017.
\item
  J. Dehaene and B. De Moor, ``Clifford group, stabilizer states, and
  linear and quadratic operations over GF(2),'' \emph{Physical Review A}
  \textbf{68}, 042318 (2003). DOI: 10.1103/PhysRevA.68.042318.
\item
  S. Aaronson and D. Gottesman, ``Improved simulation of stabilizer
  circuits,'' \emph{Physical Review A} \textbf{70}, 052328 (2004). DOI:
  10.1103/PhysRevA.70.052328.
\end{enumerate}

\end{document}
